\documentclass[11pt]{article}
\usepackage[margin=1in]{geometry}
\usepackage{amsmath,amssymb,amsthm}
\usepackage[colorlinks=true,linkcolor=blue,citecolor=blue,urlcolor=blue]{hyperref}
\hypersetup{pdftitle={Rectangular grids in crossing segment arrangements},
pdfauthor={Kaiwen Jin}}
\newtheorem{theorem}{Theorem}
\newtheorem{lemma}[theorem]{Lemma}
\newtheorem{corollary}[theorem]{Corollary}
\newcommand{\R}{\mathbb R}
\newcommand{\detv}[2]{\det(#1,#2)}
\title{Rectangular grids in crossing segment arrangements}
\author{Kaiwen Jin\\{\small\href{mailto:yui.ui.siki@gmail.com}{yui.ui.siki@gmail.com}}}
\date{}

\begin{document}
\maketitle
\begin{abstract}
For every integer \(k\geq1\), an arrangement of \(3k\) disjoint red
straight-line segments and \(3k\) disjoint blue straight-line segments,
with every red--blue pair crossing in both relative interiors, contains
\(k\) segments of each color forming an embedded rectangular grid.
This answers Question~1 of Jain, Ricci, Rollin and Schulz.
The equivalence preserves the entire subdivided segment paths, their
cyclic orders and the unbounded face.
More generally, \(n\) segments of each color contain a rectangular
subarrangement of \(p\) red and \(n-p+1\) blue segments for every
\(1\leq p\leq n\); consequently \(2k-1\) segments of each color suffice
for a \(k\)-grid.
The selection follows by averaging consecutive gaps in the circular
order of the segment directions.
\end{abstract}

Throughout, segments are closed and nondegenerate, and a crossing is
the intersection of two nonparallel supporting lines in the relative
interior of both segments. A rectangular \(m\)-by-\(q\) grid is the
arrangement of \(m\) horizontal red segments crossing \(q\) vertical
blue segments. An arrangement is equivalent to such a grid if its
subdivided colored planar graph is isomorphic to that of the grid,
preserving the counterclockwise cyclic orders at vertices and the
unbounded face. Endpoints and terminal edges are included.

In their study of the geometric thickness of \(2\)-degenerate graphs,
Jain, Ricci, Rollin and Schulz~\cite[Section~3.1, Lemma~5]{JRRS}
extract a \(k\)-grid from \(k2^{k-1}\) disjoint red segments and \(k\)
disjoint blue segments that cross completely between colors.
Their argument selects a common crossing-orientation pattern.
They ask whether \(3k\) segments of each color suffice
\cite[Section~4, Question~1]{JRRS}. We retain the orientation idea
and use the circular order of straight-segment directions to obtain
a linear bound. The geometric argument below also identifies the
endpoint tails and the unbounded face.

\begin{theorem}\label{thm:original}
Let \(k\geq1\) be an integer. Suppose that \(3k\) pairwise disjoint
red straight-line segments and \(3k\) pairwise disjoint blue
straight-line segments lie in the plane, and every red segment crosses
every blue segment in both relative interiors. Then \(k\) red segments
and \(k\) blue segments form an arrangement equivalent to a rectangular
\(k\)-by-\(k\) grid.
\end{theorem}

\section{Coherent orientations and the embedding}

An orientation of a segment chooses its initial and terminal endpoints.
Its direction vector points from the initial endpoint to the terminal
endpoint. For vectors \(u,v\in\R^2\), write
\(\detv{u}{v}=u_1v_2-u_2v_1\).

We first record the elementary quadrilateral fact used to compare
crossing orders.
\begin{lemma}\label{lem:turns}
A simple quadrilateral with four nonzero turn determinants cannot have
two positive and two negative turns. If all four turns have the same
sign, the quadrilateral is convex.
\end{lemma}
\begin{proof}
Nonzero turns exclude collinearity of any three vertices. If all four
are vertices of their convex hull, a simple cycle follows the boundary:
either of the other two cycles through the four vertices uses both intersecting
diagonals. The boundary cycle is convex and its turns have one sign.
Otherwise one vertex \(X\) lies strictly inside the triangle of the
other three. Label that triangle \(U,V,W\) counterclockwise and write
\(X=aU+bV+cW\), where \(a,b,c>0\) and \(a+b+c=1\).
Up to reversal and cyclic relabeling, the quadrilateral is \(U,X,V,W\).
Putting \(T=\detv{V-U}{W-U}>0\), its turn determinants are
\[
 bT,\qquad -cT,\qquad aT,\qquad T.
\]
There are three positive turns and one negative turn; reversal reverses
all signs. These cases prove both assertions.
\end{proof}

\begin{lemma}\label{lem:embedding}
Suppose \(m\geq1\) red segments and \(q\geq1\) blue segments are
pairwise disjoint within each color, and every red--blue pair crosses
properly. If they can be oriented with direction vectors \(u_r,v_b\)
such that
\[
 \detv{u_r}{v_b}>0
 \quad\text{for every red segment \(r\) and blue segment \(b\),}
\]
then their entire arrangement is equivalent to a rectangular
\(m\)-by-\(q\) grid.
\end{lemma}
\begin{proof}
Assume first that \(m,q\geq2\). Take two reds \(r,r'\) and two blues
\(b,b'\), and put
\[
 A=r\cap b,\quad B=r\cap b',\quad C=r'\cap b',\quad D=r'\cap b.
\]
These points are distinct. The four sides of \(ABCD\) lie on the
original segments, so nonadjacent sides are disjoint by the
same-color disjointness assumption. Adjacent sides have only their
named intersection. Thus \(ABCD\) is simple and all its turns are nonzero.

Exchange \(b,b'\) if necessary so that \(A\) precedes \(B\) on \(r\).
For some \(a,c>0\), \(\varepsilon\in\{1,-1\}\), and nonzero real
\(s,t\), write
\[
 B-A=a u_r,\qquad C-D=\varepsilon c u_{r'},\qquad
 D-A=s v_b,\qquad C-B=t v_{b'}.
\]
The four turn determinants, in order at \(A,B,C,D\), are
\begin{align*}
 \detv{A-D}{B-A}&=sa\detv{u_r}{v_b},&
 \detv{B-A}{C-B}&=at\detv{u_r}{v_{b'}},\\
 \detv{C-B}{D-C}&=\varepsilon ct\detv{u_{r'}}{v_{b'}},&
 \detv{D-C}{A-D}&=\varepsilon cs\detv{u_{r'}}{v_b}.
\end{align*}
Their signs are
\(\operatorname{sgn}s,\operatorname{sgn}t,
\varepsilon\operatorname{sgn}t,\varepsilon\operatorname{sgn}s\).
If \(\varepsilon=-1\), exactly two signs are positive and two are
negative, contradicting Lemma~\ref{lem:turns}. Hence
\(\varepsilon=1\): the two reds meet these blues in the same oriented
order. Opposite signs of \(s,t\) would again give two turns of each
sign, so \(s,t\) have the same sign. The two blues therefore meet
these reds in the same oriented order as well. All four turns agree,
and \(ABCD\) is convex.

Applying this to every pair of segments of each color shows that all
reds encounter all blues in one common oriented total order, and all
blues encounter all reds in one common oriented total order.
Label the reds \(r_1,\ldots,r_m\) in the latter order and the blues
\(b_1,\ldots,b_q\) in the former order. Map \(r_i\cap b_j\) to the
grid crossing \((j,i)\), red initial and terminal endpoints to the left
and right row endpoints, and blue initial and terminal endpoints to the
lower and upper column endpoints. These maps match exactly the
consecutive edges on every segment and preserve colors.
At a crossing, the counterclockwise rays have order
\[
 u_r,\quad v_b,\quad -u_r,\quad -v_b,
\]
since their determinant is positive. This is the order of the
rightward red and upward blue rays in the standard grid. Endpoint
vertices have degree one. Thus all cyclic orders agree.

It remains to identify the unbounded face. The four extreme segments
bound the convex quadrilateral \(Q\) with successive corners
\[
 r_1\cap b_1,\quad r_1\cap b_q,\quad
 r_m\cap b_q,\quad r_m\cap b_1.
\]
For any intermediate red, its intersections with \(b_1,b_q\) lie
strictly on the corresponding opposite sides of \(Q\), by the common
red order. Its portion between them is a chord of \(Q\), containing
all its crossings by the common blue order. The same argument applies
to intermediate blues, while the extreme portions form the sides of
\(Q\). Every crossing and every edge between crossings is therefore
in \(Q\).

The endpoints of an intermediate segment are outside \(Q\), since
its two extreme crossings are the entry and exit of its supporting
line through this convex quadrilateral, and both crossings are in
the segment's relative interior. For an extreme segment, continuation
past either corner of its side also leaves \(Q\). Outside \(Q\) the
graph consequently consists only of terminal tails, each attached to
the boundary at one end and free at the other. They have no
intersection outside \(Q\).

These tails do not disconnect the exterior from infinity. Start with
an exterior polygonal path to a sufficiently large surrounding polygon,
avoiding the finitely many endpoints and meeting tails transversely.
At a tail crossing, follow a thin neighborhood of the tail from that
crossing to its free endpoint, pass around the endpoint, and return
on the other side. This removes the crossing. The detour meets no
other tail or \(Q\): the closed tail portion in question is disjoint
from those finitely many segments and has positive clearance.
For two disjoint closed segments this clearance follows directly by
minimizing their squared distance: clamping the projection onto one
segment leaves a piecewise quadratic function on the other, whose
minimum is attained and cannot be zero. Apply this also to the four
sides of \(Q\). Repeating the finitely many detours gives a path to
infinity avoiding the graph.
Thus the exterior belongs to the unbounded face, with the extreme
cycle and all terminal tails as its boundary. All faces inside \(Q\)
are bounded because its full boundary is in the graph. The graph map
above identifies this same outer boundary in the rectangular grid.

If \(m=1\) or \(q=1\), the nontrivial crossing order is simply the
order on the sole transversal, and the graph and rotation maps still
apply. The arrangement is one segment with free half-segments attached
at its crossings. Starting with that segment, each attachment preserves
connectedness of the complement by the same free-end detour just
described. There is a unique face, the unbounded face, as in the
one-row or one-column grid. This includes \(m=q=1\).
\end{proof}

\section{Selecting segment directions}

An unoriented segment direction is an angle modulo \(\pi\), hence a
point on a circle of length \(\pi\). No red direction equals a blue
direction, since segments of those two directions could not cross
properly. Directions within one color may coincide.

\begin{lemma}\label{lem:selection}
Let \(n\geq1\), and let an arrangement have \(n\) disjoint red segments
and \(n\) disjoint blue segments, with every red--blue pair crossing
properly. For every integer \(p\) with \(1\leq p\leq n\), there are
\(p\) red and \(n-p+1\) blue segments that admit the positive
determinant orientations of Lemma~\ref{lem:embedding}.
\end{lemma}
\begin{proof}
List the red direction angles, with multiplicity, as
\(\alpha_1\leq\cdots\leq\alpha_n\) in \([0,\pi)\), and extend them by
\(\alpha_{i+n}=\alpha_i+\pi\).
Let \(g_i\) be the number of blue segment directions, counted with
multiplicity, in the open gap from
\(\alpha_i\) to \(\alpha_{i+1}\), interpreted cyclically.
Tied red directions have empty gaps. Every blue direction belongs
to exactly one gap, so
\[
 \sum_{i=1}^n g_i=n.
\]
The window of \(p\) consecutive reds starting at \(i\) spans
\([\alpha_i,\alpha_{i+p-1}]\). Its interior contains
\[
 c_i=\sum_{j=i}^{i+p-2}g_j
\]
blue directions, with the empty sum understood as zero for \(p=1\).
Every gap occurs in exactly \(p-1\) cyclic windows. Therefore
\[
 \sum_{i=1}^n c_i=(p-1)n,
\]
and some \(c_i\leq p-1\).

This chosen window has angular length strictly less than \(\pi\).
All windows have length at most \(\pi\); a window of length \(\pi\)
would contain all \(n\) blues, since no blue direction equals its
red endpoint direction, contradicting \(c_i\leq p-1<n\).
Thus at least \(n-c_i\geq n-p+1\) blues lie in its complementary
open arc.

Keep these \(p\) reds and any \(n-p+1\) such blues. Orient the reds
at their lifted angles in \([\alpha_i,\alpha_{i+p-1}]\), and the
blues at their lifts in \((\alpha_{i+p-1},\alpha_i+\pi)\).
Every blue angle exceeds every red angle by a number strictly between
zero and \(\pi\), so every red--blue determinant is positive.
The selected objects remain the original segments.
The counting also covers \(p=n\), when each window omits one gap,
and arbitrary same-color ties; no perturbation is required.
\end{proof}

\begin{proof}[Proof of Theorem~\ref{thm:original}]
Use Lemma~\ref{lem:selection} with \(n=3k\) and \(p=k\).
It selects \(k\) red and \(2k+1\) blue segments with positive
red--blue determinants. Keep any \(k\) of these blue segments.
The proper crossings, same-color disjointness and coherent orientations
are inherited. Lemma~\ref{lem:embedding} gives the required grid
with its full segment paths, cyclic orders and unbounded face.
\end{proof}

\begin{corollary}\label{cor:stronger}
Under the same disjointness and proper-crossing assumptions, \(n\)
segments of each color contain an embedded rectangular arrangement of
\(p\) red and \(n-p+1\) blue segments for every \(1\leq p\leq n\).
In particular, \(2k-1\) segments of each color suffice for a \(k\)-grid
for every integer \(k\geq1\).
\end{corollary}
\begin{proof}
Combine Lemmas~\ref{lem:selection} and~\ref{lem:embedding}; then set
\(n=2k-1\) and \(p=k\). Larger equal families can first be restricted
to \(2k-1\) of each color.
\end{proof}

Corollary~\ref{cor:stronger} is incompatible with the exclusion claimed
in \cite[Figure~26]{JRRS} for an arrangement of \(3t\) segments of each
color containing no \((t+1)\)-grid: for every \(t\geq1\),
\(3t\geq2(t+1)-1\). Under the embedded-grid definition in
Question~1 and in this note, that exclusion cannot hold.
No optimality claim is made for the sufficient bound \(2k-1\).

\begin{thebibliography}{1}
\bibitem{JRRS}
R.~Jain, M.~Ricci, J.~Rollin and A.~Schulz,
\emph{On the geometric thickness of 2-degenerate graphs},
Journal of Computational Geometry \textbf{15}(2) (2024 issue), 94--123.
Published 10 February 2025.
\href{https://doi.org/10.20382/jocg.v15i2a4}{doi:10.20382/jocg.v15i2a4}.
\end{thebibliography}

\paragraph{AI use.}
OpenAI Codex (gpt-6.1-sol) was used to develop and check the argument
and prepare the manuscript. No independent human review or formal
verification is claimed.
\end{document}
